\justifying
\NSFCBodyText

\subsubsection{已具备的实验条件}
申请人已确认可稳定使用NVIDIA A40显卡8张和NVIDIA GeForce RTX 3090 Ti显卡6张，具备Linux、Python、深度学习框架、Git及LaTeX等软件工具链。上述资源可支持公开影像处理、单卡或多卡模型训练与多种子评测，设备总量不代表本课题独占额度。MCT-P01已在单张A40上完成九组训练，纯训练计时合计约558秒，不含数据准备与验证；正式实验仍须按模型规模和可用时段测量、排期。

已取得作者公开的MineNetCD256训练与验证四分片并核验SHA256，保存了固定抽样清单、代码、模型和逐矿区指标，可用于复查MCT-P01。完整版MineNetCD归档未完整验收，官方测试影像未读取；EuroMineNet及全球铜矿数据目前仅完成来源与许可核查。正式实验继续追溯原始场景、切片位置及配准质量，按许可登记引用与衍生产物公开范围。

\subsubsection{尚缺少的实验条件及拟解决的途径}
本课题的可用算力时段、学生投入安排与经费机制为\textbf{【待补 F-MCT-05：资源排期、2026年资助金额及拨付或报销方式】}。将按确认后的资源优先保障主数据、强基线和核心消融，EuroMineNet扩展与真实人员复核依资源条件开展。研究不以购买商业影像、新增服务器或从零训练大型模型作为启动前提。

标签解释与矿山案例核验需要专业协作，拟由\textbf{【待补 F-MCT-02：遥感或地质环境合作者及核验安排】}提供支持。当前没有已落实的矿企内部监测数据、新疆矿山实测队列或专业标注团队证据。公开标签能回答的问题先独立完成，现场应用和环境风险分析另需相应观测条件。

\subsubsection{利用研究基地的计划与落实情况}
拟围绕开放课题的矿山地质环境监测方向，与中南大学有色金属成矿预测与地质环境监测教育部重点实验室开展方法交流和案例审查。具体合作人、交流及论文单位署名安排见\textbf{【待补 F-MCT-02：实验室协作与成果署名的落实情况】}。本稿未据指南推定已经获得实验室设备、账号或场地使用权；使用研究基地的具体条件以双方确认记录为准。
